% Options for packages loaded elsewhere
% Options for packages loaded elsewhere
\PassOptionsToPackage{unicode}{hyperref}
\PassOptionsToPackage{hyphens}{url}
%
\documentclass[
  11pt,
  letterpaper,
  DIV=11,
  numbers=noendperiod,
  twoside]{scrartcl}
\usepackage{xcolor}
\usepackage[left=1.1in, right=1in, top=0.8in, bottom=0.8in,
paperheight=9.5in, paperwidth=7in, includemp=TRUE, marginparwidth=0in,
marginparsep=0in]{geometry}
\usepackage{amsmath,amssymb}
\setcounter{secnumdepth}{3}
\usepackage{iftex}
\ifPDFTeX
  \usepackage[T1]{fontenc}
  \usepackage[utf8]{inputenc}
  \usepackage{textcomp} % provide euro and other symbols
\else % if luatex or xetex
  \usepackage{unicode-math} % this also loads fontspec
  \defaultfontfeatures{Scale=MatchLowercase}
  \defaultfontfeatures[\rmfamily]{Ligatures=TeX,Scale=1}
\fi
\usepackage{lmodern}
\ifPDFTeX\else
  % xetex/luatex font selection
  \setmainfont[ItalicFont=EB Garamond Italic,BoldFont=EB Garamond
Bold]{EB Garamond Math}
  \setsansfont[]{EB Garamond}
  \setmathfont[]{Garamond-Math}
\fi
% Use upquote if available, for straight quotes in verbatim environments
\IfFileExists{upquote.sty}{\usepackage{upquote}}{}
\IfFileExists{microtype.sty}{% use microtype if available
  \usepackage[]{microtype}
  \UseMicrotypeSet[protrusion]{basicmath} % disable protrusion for tt fonts
}{}
\usepackage{setspace}
% Make \paragraph and \subparagraph free-standing
\makeatletter
\ifx\paragraph\undefined\else
  \let\oldparagraph\paragraph
  \renewcommand{\paragraph}{
    \@ifstar
      \xxxParagraphStar
      \xxxParagraphNoStar
  }
  \newcommand{\xxxParagraphStar}[1]{\oldparagraph*{#1}\mbox{}}
  \newcommand{\xxxParagraphNoStar}[1]{\oldparagraph{#1}\mbox{}}
\fi
\ifx\subparagraph\undefined\else
  \let\oldsubparagraph\subparagraph
  \renewcommand{\subparagraph}{
    \@ifstar
      \xxxSubParagraphStar
      \xxxSubParagraphNoStar
  }
  \newcommand{\xxxSubParagraphStar}[1]{\oldsubparagraph*{#1}\mbox{}}
  \newcommand{\xxxSubParagraphNoStar}[1]{\oldsubparagraph{#1}\mbox{}}
\fi
\makeatother


\usepackage{longtable,booktabs,array}
\newcounter{none} % for unnumbered tables
\usepackage{calc} % for calculating minipage widths
% Correct order of tables after \paragraph or \subparagraph
\usepackage{etoolbox}
\makeatletter
\patchcmd\longtable{\par}{\if@noskipsec\mbox{}\fi\par}{}{}
\makeatother
% Allow footnotes in longtable head/foot
\IfFileExists{footnotehyper.sty}{\usepackage{footnotehyper}}{\usepackage{footnote}}
\makesavenoteenv{longtable}
\usepackage{graphicx}
\makeatletter
\newsavebox\pandoc@box
\newcommand*\pandocbounded[1]{% scales image to fit in text height/width
  \sbox\pandoc@box{#1}%
  \Gscale@div\@tempa{\textheight}{\dimexpr\ht\pandoc@box+\dp\pandoc@box\relax}%
  \Gscale@div\@tempb{\linewidth}{\wd\pandoc@box}%
  \ifdim\@tempb\p@<\@tempa\p@\let\@tempa\@tempb\fi% select the smaller of both
  \ifdim\@tempa\p@<\p@\scalebox{\@tempa}{\usebox\pandoc@box}%
  \else\usebox{\pandoc@box}%
  \fi%
}
% Set default figure placement to htbp
\def\fps@figure{htbp}
\makeatother


% definitions for citeproc citations
\NewDocumentCommand\citeproctext{}{}
\NewDocumentCommand\citeproc{mm}{%
  \begingroup\def\citeproctext{#2}\cite{#1}\endgroup}
\makeatletter
 % allow citations to break across lines
 \let\@cite@ofmt\@firstofone
 % avoid brackets around text for \cite:
 \def\@biblabel#1{}
 \def\@cite#1#2{{#1\if@tempswa , #2\fi}}
\makeatother
\newlength{\cslhangindent}
\setlength{\cslhangindent}{1.5em}
\newlength{\csllabelwidth}
\setlength{\csllabelwidth}{3em}
\newenvironment{CSLReferences}[2] % #1 hanging-indent, #2 entry-spacing
 {\begin{list}{}{%
  \setlength{\itemindent}{0pt}
  \setlength{\leftmargin}{0pt}
  \setlength{\parsep}{0pt}
  % turn on hanging indent if param 1 is 1
  \ifodd #1
   \setlength{\leftmargin}{\cslhangindent}
   \setlength{\itemindent}{-1\cslhangindent}
  \fi
  % set entry spacing
  \setlength{\itemsep}{#2\baselineskip}}}
 {\end{list}}
\usepackage{calc}
\newcommand{\CSLBlock}[1]{\hfill\break\parbox[t]{\linewidth}{\strut\ignorespaces#1\strut}}
\newcommand{\CSLLeftMargin}[1]{\parbox[t]{\csllabelwidth}{\strut#1\strut}}
\newcommand{\CSLRightInline}[1]{\parbox[t]{\linewidth - \csllabelwidth}{\strut#1\strut}}
\newcommand{\CSLIndent}[1]{\hspace{\cslhangindent}#1}



\setlength{\emergencystretch}{3em} % prevent overfull lines

\providecommand{\tightlist}{%
  \setlength{\itemsep}{0pt}\setlength{\parskip}{0pt}}



 


\KOMAoption{captions}{tableheading}
\makeatletter
\@ifpackageloaded{caption}{}{\usepackage{caption}}
\AtBeginDocument{%
\ifdefined\contentsname
  \renewcommand*\contentsname{Table of contents}
\else
  \newcommand\contentsname{Table of contents}
\fi
\ifdefined\listfigurename
  \renewcommand*\listfigurename{List of Figures}
\else
  \newcommand\listfigurename{List of Figures}
\fi
\ifdefined\listtablename
  \renewcommand*\listtablename{List of Tables}
\else
  \newcommand\listtablename{List of Tables}
\fi
\ifdefined\figurename
  \renewcommand*\figurename{Figure}
\else
  \newcommand\figurename{Figure}
\fi
\ifdefined\tablename
  \renewcommand*\tablename{Table}
\else
  \newcommand\tablename{Table}
\fi
}
\@ifpackageloaded{float}{}{\usepackage{float}}
\floatstyle{ruled}
\@ifundefined{c@chapter}{\newfloat{codelisting}{h}{lop}}{\newfloat{codelisting}{h}{lop}[chapter]}
\floatname{codelisting}{Listing}
\newcommand*\listoflistings{\listof{codelisting}{List of Listings}}
\makeatother
\makeatletter
\makeatother
\makeatletter
\@ifpackageloaded{caption}{}{\usepackage{caption}}
\@ifpackageloaded{subcaption}{}{\usepackage{subcaption}}
\makeatother
\usepackage{bookmark}
\IfFileExists{xurl.sty}{\usepackage{xurl}}{} % add URL line breaks if available
\urlstyle{same}
\hypersetup{
  pdftitle={A Question about Humberstone Frames},
  pdfauthor={Brian Weatherson},
  hidelinks,
  pdfcreator={LaTeX via pandoc}}


\title{A Question about Humberstone Frames}
\author{Brian Weatherson}
\date{2026-06-09}
\begin{document}
\maketitle
\begin{abstract}
\noindent A solution to a question Wesley Holliday asked about
Humberstone frames.
\end{abstract}


\setstretch{1.1}
This post concerns \textbf{Problem 8.13} in Holliday
(\citeproc{ref-Holliday2025}{2025}), which asks, roughly, whether the
argument of §2.5 of that paper can be extended to what he calls
\emph{Humberstone frames}. The answer is yes, as I'll show here.

This is a blog post, not a paper, so I won't repeat any definitions that
are already in Holliday's paper, or for that matter in Humberstone
(\citeproc{ref-Humberstone1981a}{1981}). The only thing I will note is
that I'm following Humberstone, and not Holliday, in using \(x \leq y\)
to mean that \(y\) is a refinement of \(x\). Relatedly, \(x \geq y\)
just means \(y \leq x\).

The frame we'll use consists of (a) all finite binary strings, plus (b)
a `code' for each binary string. When \(y\) is a code, I'll write
\([y]\) for the string that \(y\) is the code of, and when \(x\) is a
string, I'll write \(c(x)\) for the code for \(x\). The refinement
relation \(x \leq y\) holds in one of two cases:

\begin{enumerate}
\def\labelenumi{\arabic{enumi}.}
\tightlist
\item
  \(x=y\)
\item
  \(x\) and \(y\) are strings, and \(x\) is an initial segment of \(y\).
\end{enumerate}

That is, strings are refined by adding to them, and codes have no
non-trivial refinements. We then introduce a bunch of accessibility
relations: \(R_i\), which determine modal operators \(\Box_i\), and
hence by definition \(\diamond_i\).

\begin{description}
\tightlist
\item[Universal \(R_\forall\)]
\(xR_{\forall}y\) whenever \(x\) and \(y\) are points in the model.
\item[String-to-code \(R_{\rightarrow}\)]
\(xR_{\rightarrow}y\) iff \(x\) is a string, \(y\) is a code, and one of
\(x\) and \([y]\) is an initial segment of the other. I'm imagining the
strings on the left, and the codes on the right, hence the left-to-right
arrow.
\item[Code-to-string \(R_{\leftarrow}\)]
\(xR_{\leftarrow}y\) iff \(x\) is a code, \(y\) is a string, and \([x]\)
is an initial segment of \(y\). Note this is not the same definition as
for \(R_{\rightarrow}\), since here we insist on which is the initial
segment.
\item[Code-zero-remove: \(R_{-}\)]
\(xR_{-}y\) iff \(x\) and \(y\) are codes, and \([x]\) ends with a 0,
and \([y]\) is \([x]\) minus that trailing 0.
\item[Code-shortening: \(R_{\subsetneq}\)]
\(xR_{\subsetneq}y\) iff \(x\) and \(y\) are distinct codes, and \([y]\)
is an initial segment of \([x]\)
\end{description}

Call the following proposition \(\alpha\)

\[
\diamond_{-}(\Box_{\leftarrow} p \wedge \neg \diamond_{\subsetneq}\Box_{\leftarrow} p)
\]

Let's build up where this will be true. (Note that I'm following
Holliday's own construction \emph{really closely} here.)

\begin{itemize}
\tightlist
\item
  \(\Box_{\leftarrow} p\) is true at the code \(x\) iff \(p\) is true at
  \([x]\).
\item
  So \(\diamond_{\subsetneq}\Box_{\leftarrow} p\) is true at \(x\) if
  \(p\) is true at some proper initial segment of \([x]\).
\item
  So \(\neg \diamond_{\subsetneq}\Box_{\leftarrow} p\) is true at \(x\)
  if \(p\) is not true at any proper initial segment of \([x]\).
\item
  So
  \(\Box_{\leftarrow} p \wedge \neg \diamond_{\subsetneq}\Box_{\leftarrow} p\)
  is true at \(x\) if \(p\) is true at \([x]\), but not at any
  shortening of \([x]\). That is, if you were working down the tree of
  strings, \([x]\) is the first point you'd get to where \(p\) is true.
  Call a string like this an \emph{initial point} for \(p\).
\item
  So
  \(\diamond_{-}(\Box_{\leftarrow} p \wedge \neg \diamond_{\subsetneq}\Box_{\leftarrow} p)\)
  is true at the code \(x\) if \([x]\) is \(y\frown \langle0\rangle\),
  where \(y\) is an \emph{initial point} for \(p\).
\end{itemize}

Now consider any initial point \(x\) for \(p\). Since \(p\) is
persistent, \(p\) will be true at \(x\frown \langle0\rangle\), call this
\(x_0\), and at \(x\frown \langle1\rangle\), call this \(x_1\). Since,
as we just showed, \(\alpha\) is true at \(c(x_0)\), and
\(x_0R_{\rightarrow}c(x_0)\), it follows that
\(\diamond_{\rightarrow} \alpha\) is true at \(x_0\). But
\(\diamond_{\rightarrow} \alpha\) will not be true at \(x_1\), or at any
of its refinements. At \(c(x_1)\), \(\alpha\) fails because this is not
the code for a sequence ending in 0. And at the code for any shortening
or extension of \(x_1\), \(\alpha\) fails because that's not the code of
a string one element longer than an \emph{initial} point for \(p\).

That means that if \(p\) is true at some string, it is true at some
initial point, and hence both \(p \wedge \diamond_{\rightarrow} \alpha\)
and \(p \wedge \neg \diamond_{\rightarrow} \alpha\). If
\(\diamond_{\forall}(p \wedge \diamond_{\rightarrow} \top)\) is true
(anywhere), then \(p\) is true at some string. (It can't be that \(p\)
is simply true at codes, because \(\diamond_{\rightarrow} \top\) is not
true at any code.) So
\(\diamond_{\forall}(p \wedge \diamond_{\rightarrow} \alpha)\) and
\(\diamond_{\forall}(p \wedge \neg \diamond_{\rightarrow} \alpha)\) will
both be true everywhere. Hence at all points, we'll have this instance
of (Split):

\begin{equation}\protect\phantomsection\label{eq-split}{
\diamond_{\forall}(p \wedge \diamond_{\rightarrow} \top) \rightarrow [\diamond_{\forall}(p \wedge \diamond_{\rightarrow} \alpha) \wedge \diamond_{\forall}(p \wedge \neg \diamond_{\rightarrow} \alpha)] 
}\end{equation}

It just remains to show that all the \(R\)-relations satisfy all the
conditions for Humberstone-frames. This is trivial for \(R_\forall\),
because it always holds, and for \(R_{\subsetneq}\) and \(R_{-}\)
because they just relate points that don't stand in any non-trivial
refinement relations. The only questions are for \(R_{\rightarrow}\) and
\(R_{\leftarrow}\). Let's go through the three-conditions on
Humberstone-frames, using Holliday's names for them.

First \textbf{upR}, i.e.,
\(x \leq x^{\prime} \wedge x^{\prime}Ry \rightarrow xRy\). This can't
fail for \(R_{\leftarrow}\), because the left-hand side has no
non-trivial coarsenings. For \(R_{\rightarrow}\),
\(x^{\prime}R_{\rightarrow}y\) means either \(x^\prime\) is an initial
segment of \([y]\), or \([y]\) is an initial segment of \(x^\prime\). If
the former, then clearly \(x\), which is an initial segment of
\(x^\prime\), is also an initial segment of \([y]\), so
\(xR_{\rightarrow}y\). If the latter, imagine working through all the
coarsenings, i.e., shortenings, of \(x^\prime\) sequentially by starting
with \(x^\prime\) and removing one element at a time. At first we'll
have sequences that are, like \(x^\prime\), extensions of \([y]\), then
we'll get to \([y]\) itself, then we'll have shortenings of \([y]\). At
every stage, we'll have something that is either an extension of
\([y]\), or a shortening of it, so \(xR_{\rightarrow}y\) will hold.

Now \textbf{Rdown}, i.e.,
\(xRy \wedge y \leq y^\prime \rightarrow xRy^\prime\). In this case,
this is trivial for \(R_{\rightarrow}\), since when
\(xR_{\rightarrow}y\), \(y\) has no non-trivial refinements. For
\(R_{\leftarrow}\), \(xR_{\leftarrow}y\) means that \(y\) is an
extension of \([x]\), so any extension of \(y\) will be an extension of
\([x]\), as required.

Now \textbf{R-ref++}, i.e.,
\(xRy \rightarrow \exists x^\prime \geq x \forall x^{\prime\prime} \geq x^\prime: x^{\prime\prime}Ry\).
Again, this is trivial for \(R_{\leftarrow}\), since whenever
\(xR_{\leftarrow}y\), \(x\) has no non-trivial refinements. The only
tricky case is \(R_{\rightarrow}\). So assume \(xR_{\rightarrow}y\).
Again there are two cases to consider. One is where \(x\) is an initial
segment of \([y]\). In that case it we can set \(x^\prime\) to be
\([y]\). The other is where \([y]\) is an initial segment of \(x\). In
that case, we can just set \(x^\prime\) to be \(x\) itself. Either way,
there is a witness for the consequent of \textbf{R-ref++}.

The result is that Equation~\ref{eq-split} will be valid on the frame.
Here we follow Holliday to complete the proof. There is no class of
Kripke frames where Equation~\ref{eq-split} is valid on the frames
without
\(\neg \diamond_{\forall}(p \wedge \diamond_{\rightarrow} \top)\) being
valid on the frames. Quick proof: if that were not the case, so there
was a frame where
\(\diamond_{\forall}(p \wedge \diamond_{\rightarrow} \top)\) could be
true, there would be a model on that frame where \(p\) was only true at
one point, so there could not be witnesses for both
\(\diamond_{\forall}(p \wedge \diamond_{\rightarrow} \alpha)\) and
\(\diamond_{\forall}(p \wedge \neg \diamond_{\rightarrow} \alpha)\). So
there is a class of Humberstone frames (in fact a singleton class),
which determines a logic not determinable by any class of Kripke frames.

\section*{References}\label{references}
\addcontentsline{toc}{section}{References}

\protect\phantomsection\label{refs}
\begin{CSLReferences}{1}{0}
\bibitem[\citeproctext]{ref-Holliday2025}
Holliday, Wesley H. 2025. {``Possibility Frames and Forcing for Modal
Logic.''} \emph{Australasian Journal of Logic} 22 (2): 44--288. doi:
\href{https://doi.org/10.26686/ajl.v22i2.5680}{10.26686/ajl.v22i2.5680}.

\bibitem[\citeproctext]{ref-Humberstone1981a}
Humberstone, Lloyd. 1981. {``From Worlds to Possibilities.''}
\emph{Journal of Philosophical Logic} 10 (3): 313--39. doi:
\href{https://doi.org/10.1007/BF00293423}{10.1007/BF00293423}.

\end{CSLReferences}




\end{document}
